\begin{recipe}{Gnocchi}
\label{recipe:gnocchi}
\kicker[Hoofdgerecht,Italiaans,Basisrecept]{Hoofdgerecht · Italiaans · basisrecept}
\index[register]{Gnocchi}
\index[register]{Aardappelen}
\index[register]{Bloem}

\meta{Ca. 1,5 uur (vooral wachttijd)}

\lettrine{Z}{elfgemaakte} gnocchi vraagt vooral geduld: de aardappelpuree
moet goed afgekoeld zijn voor je de bloem erdoor mengt, anders wordt het
deeg plakkerig. Gebruik niet meer bloem dan nodig, dat maakt de gnocchi
taai.

\blockrule{Ingrediënten}
\begin{ingredients}
  \ing{1 kg}{kruimige aardappelen}\index[register]{Aardappelen}
  \ing{250--300 g}{bloem type 00 (tipo 00)}\index[register]{Bloem}
  \ing{1}{ei}
  \ing{1 tl}{zout}
  \ingb{extra bloem, om te bestuiven}
\end{ingredients}

\blockrule{Bereiding}
\tip{Gebruik een pureeknijper (wij gebruiken de Joseph Joseph Helix) in
plaats van een staafmixer. Een staafmixer maakt de aardappel te plakkerig
voor gnocchi.}
\begin{steps}
  \item Kook de aardappelen met schil in ruim water gaar, 20--30 minuten.
  \item Giet ze af, laat ze kort uitdampen en pel ze terwijl ze nog warm
        zijn.
  \item Druk de aardappelen fijn met een pureeknijper. Laat de puree
        volledig afkoelen.
  \item Meng het ei en het zout kort door de afgekoelde puree.
  \item Voeg geleidelijk de bloem toe en kneed tot een zacht, soepel deeg.
        Kneed niet te lang.
  \item Verdeel het deeg in stukken en rol elk stuk uit tot een rol van
        ongeveer 2 cm dik.
  \item Snijd de rollen in stukjes van ongeveer 2 cm.
  \item Rol elk stukje eventueel over de achterkant van een vork of een
        gnocchiplankje voor de kenmerkende ribbels, en druk met je duim een
        kuiltje in het midden.
  \item Breng een grote pan gezouten water aan de kook en kook de gnocchi in
        porties. Zodra ze na 2--3 minuten bovendrijven, zijn ze gaar.
\end{steps}
\end{recipe}
